\begin{figure*}[!tb]
\centering
\begin{tikzpicture}[
  font=\footnotesize,
  % Rounded = a processing step, square-cornered = a data object that flows between steps.
  bx/.style={draw, minimum height=7mm, inner xsep=4pt, align=center},
  proc/.style={bx, rounded corners=2pt},
  obj/.style={bx, sharp corners},
  data/.style={obj, fill=cInput!18, draw=cInput!70!black},
  objmid/.style={obj, fill=black!4, draw=black!45},
  learn/.style={proc, fill=cTrain!18, draw=cTrain!70!black},
  flowb/.style={proc, fill=black!8, draw=black!55},
  ar/.style={-{Latex[length=2mm]}, draw=black!70},
  lbl/.style={anchor=north west, text width=0.97\textwidth, align=left},
]

% ---------- what the stage has to undo: the flicker itself ----------
% Sits well clear of panel (a): at a smaller yshift the sketch's axis ran through the
% "point tracking" box and its y-label collided with the orange note.
\begin{scope}[shift={(0.35,2.75)}]
  \draw[->, black!55] (0,-0.5) -- (4.5,-0.5) node[right, black!70, font=\scriptsize] {frame};
  \draw[->, black!55] (0,-0.56) -- (0,0.66);
  \node[rotate=90, anchor=south, font=\scriptsize, text=black!70] at (-0.42,0.05)
        {depth scale};
  \draw[dashed, black!45] (0,0) -- (4.25,0)
        node[right, black!70, font=\scriptsize, inner xsep=4pt] {true};
  \draw[thick, cInput!75!black]
    (0.15,0.26) -- (0.7,-0.20) -- (1.25,0.32) -- (1.8,-0.09) -- (2.35,0.38)
    -- (2.9,-0.28) -- (3.45,0.15) -- (4.0,-0.22);
  \node[anchor=west, font=\scriptsize, text=cInput!60!black] at (-0.1,1.02)
        {DA3-g is off by a different factor every frame: still points seem to breathe};
  \node[anchor=west, font=\scriptsize, align=left] at (5.6,0.0)
        {The stage emits \emph{one} number per frame, $\Delta s_f$, and scales\\
         that frame's depth by $e^{\Delta s_f}$ to flatten the wobble onto\\
         the dashed line. That drift comes from the depth model alone.};
\end{scope}

% ---------------- (a) the de-flicker, as built ----------------
% Boxes carry process names only. The earlier version spelled the contents out inside the nodes,
% which made the "tracked points" box wide enough to reach the pooling box and left the elbow
% arrows running along its edge; the detail now lives in the caption.
\node[data]  (aimg)  at (0,0)      {video\\frames};
\node[flowb] (aflow) at (2.2,0)    {optical\\flow};
\node[data]  (adep)  at (0,-1.6)   {DA3-g\\depth map};
\node[flowb] (apts)  at (4.5,0)    {point\\tracking};
\node[flowb] (asamp) at (4.5,-1.6) {depth\\sampling};
\node[learn] (atok) at (8.0,-0.8)  {masked mean\\$\rightarrow$ 1 token/frame};
\node[learn] (amam) at (11.0,-0.8) {Mamba-3\\over frames};
\node[objmid] (aout) at (13.3,-0.8) {$\Delta s_f$};

\draw[ar] (aimg) -- (aflow);
\draw[ar] (aflow) -- (apts);
\draw[ar] (apts) -- (asamp);
\draw[ar] (adep) -- (asamp);
% One shared elbow into the pooling box, clear of both source boxes.
\draw[ar] (apts.east)  -- ++(0.55,0) |- (atok.west);
\draw[ar] (asamp.east) -- ++(0.55,0) |- (atok.west);
\draw[ar] (atok) -- (amam);
\draw[ar] (amam) -- (aout);

\node[lbl] at (0,-2.35)
  {\textbf{(a) De-flicker stage, as built.} \emph{Point tracking} propagates the clip's
   annotated query points --- $806$--$901$ on adt, $256$ on drivetrack, as few as $98$ on
   pstudio --- to every frame with the frozen optical flow, and marks each
   one visible or not by a forward--backward consistency check. \emph{Depth sampling} reads the
   DA3-g depth map at exactly those pixels and nowhere else. The \emph{masked mean} then collapses
   the points to one vector per frame, weighting each by its visibility, so Mamba-3 sees a sequence
   of $F$ frame tokens and never sees a point. Everything the stage observes therefore arrives
   \emph{through the tracker}: a stage trained with SEA-RAFT and one trained with WAFT learn
   different corrections from identical depth --- even though the wobble they must undo belongs to
   the depth model alone and has nothing to do with tracking.};

% ---------------- (b) the Mamba-3 depth scale refiner ----------------
\begin{scope}[yshift=-6.6cm]
\node[data]  (bdep)  at (0,-0.8)    {DA3-g\\depth map};
\node[flowb] (bds) at (2.9,-0.8) {take log,\\pool to $64^2$};
\node[learn] (bcnn)  at (5.9,-0.2)  {small CNN\\$\rightarrow$ 64 features};
\node[flowb] (bq) at (5.9,-1.5) {10 log-depth\\quantiles};
\node[learn] (btok) at (8.9,-0.8)   {concat\\$\rightarrow$ 1 token/frame};
\node[learn] (bmam) at (11.6,-0.8)  {Mamba-3\\over frames};
\node[objmid] (bout) at (13.7,-0.8)  {$\Delta s_f$};

\draw[ar] (bdep) -- (bds);
\draw[ar] (bds.east) -- ++(0.3,0) |- (bcnn.west);
\draw[ar] (bds.east) -- ++(0.3,0) |- (bq.west);
\draw[ar] (bcnn.east) -- ++(0.3,0) |- (btok.west);
\draw[ar] (bq.east)   -- ++(0.3,0) |- (btok.west);
\draw[ar] (btok) -- (bmam);
\draw[ar] (bmam) -- (bout);

\node[lbl] at (0,-2.15)
  {\textbf{(b) Mamba-3 depth scale refiner.} It reads the depth map itself, so no tracker appears
   anywhere in the path and the front-end cannot change what it learns, and it sees every pixel of
   the frame rather than the few hundred scattered points panel (a) is restricted to. Note what the
   CNN emits: \emph{$64$ features}, not a number. Those are concatenated with ten quantiles of the
   frame's log depth --- the CNN ends in a global average, which discards the shape of the depth
   distribution, and the quantiles put it back --- and projected to one $128$-dimensional token per
   frame. Only after Mamba-3 has mixed the $F$ tokens does a linear head produce the scalar
   $\Delta s_f$. Both stages therefore hand Mamba-3 the same thing: one token per frame.};
\end{scope}

\end{tikzpicture}
\caption{\textbf{What the de-flicker stage is for, and why the way it was built was wrong.}
  \emph{The problem it solves:} DA3-g estimates each frame independently, so its depth is off by a
  slightly different factor every frame --- the sketch at the top. Points that are still in the
  world therefore appear to breathe toward and away from the camera. The stage exists to undo that,
  and emits one number per frame, $\Delta s_f$, that rescales that frame's depth.
  \emph{Why (a) was a problem:} the drift is a property of the depth model, but as built the stage
  can only observe it through the tracker --- depth sampled at flow-chosen pixels, averaged under a
  flow-derived visibility mask. That is what couples it to the front-end, and why swapping SEA-RAFT
  for WAFT changed its behaviour on identical depth. (b) removes the coupling by reading the map.
  \emph{Why it also learned little:} as built there is no de-flicker loss at all. The stage is
  trained only through the tracking loss
  $L=vis\cdot|xyz_{\text{pred}}-xyz_{\text{gt}}|_1/\text{scale}_{\text{gt}}+\lambda_{uv}vis|\Delta
  uv|^2$, which sums over the clip's $98$--$901$ points $\times$ $F$ frames, while the stage controls a single
  per-frame scalar worth only $0.3$--$0.6\%$ of that error (\S\ref{sec:deflicker-signal}): about
  $99.5\%$ of its gradient is noise from its point of view. Giving it a direct target,
  $\Delta s^{*}_f=\mathrm{median}_n[\log z_{\text{gt}}-\log z_{\text{raw}}]$ with
  $L_{\text{dsr}}=|\Delta s_f-\Delta s^{*}_f|$, supervises exactly what it can change.
  Orange is data, blue is trainable, grey is frozen or intermediate; rounded boxes are
  processing steps and square-cornered boxes are the data objects passed between them.}
\label{fig:deflicker-coupling}
\end{figure*}
